% Gabriel Ryan's Resume

%\documentclass[9pt,letterpaper]{article} % Font size and paper type
\documentclass{article} % Font size and paper type

\usepackage[parfill]{parskip} % Remove paragraph indentation
\usepackage{array} % Required for boldface (\bf and \bfseries) tabular columns
\usepackage[hidelinks]{hyperref} % Required for boldface (\bf and \bfseries) tabular columns
%\usepackage{ifthen} % Required for ifthenelse statements
\usepackage{geometry}
\usepackage{color}
\usepackage[usenames,dvipsnames]{xcolor}
\usepackage{graphicx}
\usepackage{verbatim} %multiline comments
\usepackage{url}
\usepackage{enumitem}

\definecolor{lightgrey}{gray}{.95}


%\usepackage{fontspec}
\usepackage[sfdefault,scaled=0.92]{noto-sans}
\usepackage[scale=0.90]{noto-mono}

%\usepackage[scaled]{berasans}
%\renewcommand*\familydefault{\sfdefault}  %% Only if the base font of the document is to be sans serif

%\usepackage[usefilenames,RMstyle=Light,SSstyle=Light,TTstyle=Light,DefaultFeatures={Ligatures=Common}]{plex-otf} %
%\renewcommand*\sfdefault{lmssq}
%\usepackage[sfdefault]{roboto}  %% Option 'sfdefault' only if the base font of the document is to be sans serif
\renewcommand*\familydefault{\sfdefault} %% Only if the base font of the document is to be sans serif
%\usepackage[T1]{fontenc}
%\defaultfontfeatures{Mapping=tex-text,Scale=MatchLowercase}

%\usepackage{helvet}
%\usepackage{lmodern}
%\usepackage{tgtermes}
%\usepackage{palatino}


\geometry{right=1in,left=.8in,top=.6in,bottom=.7in}
\pagestyle{empty} % Suppress page numbers

\newcommand{\tbf}[1]{\textbf{#1}}

%%%%%%%%%%%%%%%%%%%%%%%%% Commands %%%%%%%%%%%%%%%%%%%%%%%%%
%\rsection{<section title>}
\newcommand{\rsection}[1]{
  \hspace{-0.4cm}\vspace{0.1cm}
\colorbox{lightgrey}{
\begin{minipage}{1.07\linewidth}
\vspace{0.22cm}
\fontsize{14pt}{16pt}\selectfont #1
%\rule[9pt]{\linewidth}{1.2pt}
\vspace{0.12cm}
\end{minipage}
}
\vspace*{-0.1cm}
}
%\rjob{<job title>}{<time>}
\newcommand{\rjob}[2]{
  \hspace*{-0.3cm}
{\fontsize{10pt}{12pt}\selectfont #1} \hfill #2
\hspace*{-1.2cm}
}
%\begin{ritemize}...\end{ritemize}
\newenvironment{ritemize}{
\hspace*{-0.8cm}
\begin{minipage}{1.05\linewidth}
\vspace{0.1cm}
\begin{itemize}[itemsep=0pt,parsep=0pt,topsep=2pt,partopsep=0pt]
}{
\end{itemize}
\end{minipage}
%\vspace{-0.1cm}
}
%\ritem
\newcommand{\ritem}{
\item[-]
%\hspace{-8pt}
}


\begin{document}
%%%%%%%%%%%%%%%%%%%%%%%%% Header %%%%%%%%%%%%%%%%%%%%%%%%%
%\begin{center}
\hspace*{-0.45cm}
{\fontsize{22pt}{22pt}\selectfont Gabriel Ryan}\\
%\vspace{1.2cm}
%\vspace*{-0.5cm}
%\hspace*{0.1cm}
\begin{minipage}{\linewidth}
\vspace{0.1cm}
  {\fontsize{11}{13}\selectfont
      \href{mailto:gabriel.ryan11@gmail.com}{gabriel.ryan11@gmail.com} \ \textbar\ \href{https://linkedin.com/in/gabrielryan1}{LinkedIn} \ \textbar\ \href{https://gryan11.github.io}{Website} \ \textbar\ \href{https://scholar.google.com/citations?user=6mW4eCcAAAAJ\&hl=en}{Google Scholar}
  }
\end{minipage}
\vspace{-0.15cm}

%%%%%%%%%%%%%%%%%%%%%%%%% Research Summary %%%%%%%%%%%%%%%%%%%%%%%%%
\rsection{Research Summary}

\begin{minipage}{\linewidth}
Research scientist working on AI-driven software development. My work spans model training and harness design, including completion, context retrieval, and agentic code review now shipping in GitHub Copilot to millions of developers. Before that, my PhD research combined program analysis and machine learning, developing neural methods for program verification, test generation, and vulnerability detection. I'm interested in how far agents can be pushed on real software engineering tasks, particularly post-training and multi-agent design.
\end{minipage}
\vspace{0.15cm}

%%%%%%%%%%%%%%%%%%%%%%%%% Education %%%%%%%%%%%%%%%%%%%%%%%%%
\rsection{Education}

\rjob{\textbf{Ph.D., Computer Science}, Columbia University \textit{(advised by Prof. Suman Jana)}}{Dec 2023}\\
\rjob{\textbf{M.S., Computer Science}, Columbia University}{Feb 2018}\\
\rjob{\textbf{B.S. Engineering \& B.A. Computer Science}, Swarthmore College}{Jun 2013}\\
\vspace{-0.05cm}



%%%%%%%%%%%%%%%%%%%%%%%%% Experience %%%%%%%%%%%%%%%%%%%%%%%%%
\rsection{Experience}

\rjob{\textbf{Senior Research Scientist}, Microsoft CodeAI}{Mar 2024 -- Present}\\
\begin{ritemize}
    \ritem Train custom models for Copilot subagents, context compaction, and memory generation using the Slime RL framework. Achieved a 15\% cost reduction with custom subagent models in internal deployments.
    \ritem Led development of an efficient code embedding model used to index GitHub for Copilot's context retrieval, reducing indexing cost by $2\times$ and index size by $3\times$ while improving retrieval quality.
    \ritem Trained the code completion models powering GitHub Copilot in editors used by millions of developers, driving a 5\% increase in acceptance rate while reducing latency by 10\% through custom code datasets filtered for quality and relevance.
\end{ritemize}

\rjob{\textbf{Doctoral Researcher}, Columbia Security Lab}{Sep 2018 -- Dec 2023}\\
\begin{ritemize}
    \ritem Built a novel Linux kernel race prediction analysis using probabilistic and spectral learning. Published in Oakland S\&P 2023.
    \ritem Designed \textit{Proximal Gradient Analysis} (PGA), a dataflow analysis method based on nonsmooth optimization. Implemented in LLVM. Published in USENIX Security 2021.
    \ritem Designed the \textit{Continuous Logic Network} (CLN), a neural architecture for logical learning and reasoning; applied it to learning program invariants for verification. Released a PyTorch library for constructing CLNs. Published in ICLR 2020 and PLDI 2020.
\end{ritemize}

\rjob{\textbf{Research Intern}, AWS AI Labs}{Jun 2023 -- Aug 2023}\\
\begin{ritemize}
  \ritem Built a novel LLM regression-testing approach using static analysis to prompt the model to reason symbolically about program execution paths. Achieved $2\times$ coverage and $3\times$ correct test generations over baseline approaches on CodeGen2, GPT-3.5, and GPT-4.
  \ritem Published \textit{Code-Aware Prompting} at FSE 2024.
\end{ritemize}


\rjob{\textbf{Research Intern}, Microsoft Research}{Jun 2021 -- Aug 2021}\\
\begin{ritemize}
    \ritem Built \textit{TOGA: A Neural Method for Test Oracle Generation} using Transformers (CodeBERT) and a specialized grammar to generate unit tests, demonstrating a 170\% improvement over prior work in bug detection.
    \ritem Published at ICSE 2022; awarded ACM SIGSOFT Distinguished Paper Award.
\end{ritemize}

\rjob{\textbf{Software Engineer}, Allure Security Technology}{Sep 2015 -- Aug 2018}\\
\begin{ritemize}
\ritem Built systems for managing large-scale deployments of user-behavior sensors and detecting and tracking data loss across enterprise networks.
\end{ritemize}

\rjob{\textbf{Robotics Software Consultant}, 3DDataLtd}{May 2015 -- Aug 2015}\\
\begin{ritemize}
\ritem Built a LiDAR- and inertial-sensor-based 3D mapping framework for robotic navigation.
\end{ritemize}

\rjob{\textbf{Radar Software Engineer}, Raytheon}{Sep 2013 -- Apr 2015}\\
\begin{ritemize}
\ritem Built radar calibration, satellite tracking, and antenna diagnostic software; earned a team achievement award for delivering new diagnostic capabilities ahead of schedule.
\end{ritemize}



%\clearpage{}



%\rjob{\textbf{Writing Associate}, Swarthmore College}{Sep 2011- Dec 2012}\\
%\begin{ritemize}
%\ritem Selected for excellent written and oral communication skills to work in Swarthmore writing program.
%\ritem Conducted half hour conferences with students for revision and good writing practice in college writing center.
%\ritem Assisted one professor per semester in teaching a writing intensive course, meeting regularly with 10-14 students to work on writing skills and assist in planning and revising papers.
%\end{ritemize}


%\rjob{\textbf{IT Associate}, Swarthmore College}{Fall 2011}\\
%\begin{ritemize}
%\ritem Selected for excellent communication skills and practical problem solving ability to work on college IT team.
%\ritem Fielded calls and walk-ins at IT support office. Addressed security and networking issues on PCs.
%\end{ritemize}



%\rjob{\textbf{Intern}, Verizon Laboratories}{January 2011}\\
%\begin{ritemize}
%\ritem Shadowed Verizon Interface Design Engineer developing Videochat interface for smartphones.
%\ritem Evaluated usability of current videochat apps and prepared recommendations for videochat interface design.
%\end{ritemize}

%%%%%%%%%%%%%%%%%%%%%%%%% Extracurricular %%%%%%%%%%%%%

\vspace{0.2cm}
\rsection{Selected Publications}

\hspace*{-0.1cm}
\begin{minipage}{1.01\linewidth}
\begin{itemize}[label={},itemindent=-2em,leftmargin=2em, parsep=5pt]
  \item {\bf [ICML 2026]}
    \textit{SemRep: Generative Code Representation Learning with Code Transformations.}
    W. Li, J. Song, B. A. Stoica, A. Dhoot, \textbf{G. Ryan}, S. Fu, K. Pei.
  \item {\bf [NeurIPS DL4C Workshop 2025]}
    \textit{DevBench: A Realistic, Developer-Informed Benchmark for Code Generation Models.}
    P. A. Golnari, A. Kumarappan, W. Wen, X. Liu, \textbf{G. Ryan}, Y. Sun, S. Fu, et al.
  \item {\bf [OOPSLA 2024]}
  \textit{Accurate Data Race Prediction in the Linux Kernel through Sparse Fourier Learning.} 
    \textbf{G. Ryan}, B. Cetin, Y. Lim, S. Jana. \url{https://dl.acm.org/doi/pdf/10.1145/3649840}
  \item {\bf [FSE 2024]}
    \textit{Code-Aware Prompting: A study of Coverage Guided Test Generation in Regression Setting using LLM.} 
      \textbf{G. Ryan}, S. Jain, M. Shang, S. Wang, X. Ma, M. Ramanathan, B. Ray. \url{https://arxiv.org/pdf/2402.00097}
  \item {\bf [Oakland S\&P 2023]}
  \textit{Precise Detection of Kernel Data Races with Probabilistic Lockset Analysis.}
    \textbf{G. Ryan}, A. Shah, D. She, S. Jana. \url{https://cs.columbia.edu/~gabe/files/oakland2023_pla.pdf}
  \item {\bf [ICSE 2022]}
    %\textrm{\fontsize{11pt}{12pt}\selectfont TOGA: A Neural Method for Test Oracle Generation.}
    \textit{TOGA: A Neural Method for Test Oracle Generation.}
    E. Dinella*, \textbf{G. Ryan}*, T. Mytkowicz, S. Lahiri. \url{https://arxiv.org/pdf/2109.09262.pdf} \textbf{(Distinguished Paper Award)}
  \item {\bf [OSDI 2021]}
    \textit{DistAI: Data-Driven Automated Invariant Learning for Distributed Protocols.}
    J. Yao, R. Tao, R. Gu, J. Nieh, S. Jana, \textbf{G. Ryan}. \url{https://www.usenix.org/system/files/osdi21-yao.pdf} \textbf{(Best Paper Award)}
  \item {\bf [USENIX Security 2021]}
    \textit{Fine Grained Dataflow Tracking with Proximal Gradients.}
    \textbf{G. Ryan}, A. Shah, D. She, K. Bhat, and S. Jana. \url{https://arxiv.org/abs/1909.03461}
  \item {\bf [PLDI 2020]}
    \textit{Learning Nonlinear Loop Invariants with Gated Continuous Logic Networks.}
    J. Yao*, \textbf{G. Ryan}*, J. Wong*, S. Jana, and R. Gu. \url{https://arxiv.org/abs/2003.07959}
  \item {\bf [ICLR 2020]}
    \textit{CLN2INV: Learning Loop Invariants with Continuous Logic Networks.}
    \textbf{G. Ryan}*, J. Wong*, J. Yao*, R. Gu, and S. Jana. \url{https://arxiv.org/abs/1909.11542}
\end{itemize}
\end{minipage}

\vspace{-0.1cm}
{\footnotesize * denotes equal contribution.}

\vspace{0.05cm}
\rsection{Awards}
\hspace*{-0.1cm}

\begin{minipage}{1.01\linewidth}
\begin{itemize}[label={},itemindent=-2em,leftmargin=2em,parsep=1pt,itemsep=0pt,topsep=3pt]
  \item {\bf ACM SIGSOFT Distinguished Paper Award}, ICSE 2022 (TOGA).
  \item {\bf Jay Lepreau Best Paper Award}, OSDI 2021 (DistAI).
  \item {\bf National Defense Science and Engineering Graduate Fellowship} (NDSEG). Awarded NDSEG Fellowship for proposal "Proximal Gradient Analysis for Vulnerability Detection and Defense." 2019.
  \item {\bf NSF Graduate Research Fellowship Honorable Mention}. Accorded honorable mention for proposal "Modeling and Simulating Adversarial User Behavior with Sequential VAEs." 2018.
  \end{itemize}
\end{minipage}

\vspace{0.05cm}

\rsection{Academic Service}
\hspace*{-0.1cm}
\begin{minipage}{1.01\linewidth}
\begin{itemize}[label={},itemindent=-2em,leftmargin=2em,parsep=1pt,itemsep=0pt,topsep=3pt]
  \item Workshop on Large Language Models for Code, ICSE 2026. Program Committee.
  \item Next-Gen Code Development with Collaborative AI Agents Workshop, AAAI 2026. Program Committee.
  \item AAAI 2024, 2025, 2026. Program Committee.
  \item TOSEM 2023, 2024. Reviewer.
  \item AAAI 2023. Reviewer.
  \end{itemize}
\end{minipage}

%%%%%%%%%%%%%%%%%%%%%%%%% Technical Skills %%%%%%%%%%%%%%%%%%%%%%%%%
\vspace{0.05cm}
\rsection{Technical Skills}
\hspace*{-0.1cm}
\begin{minipage}{1.01\linewidth}
\begin{itemize}[label={},itemindent=-2em,leftmargin=2em,parsep=1pt,itemsep=0pt,topsep=3pt]
  \item \textbf{Languages:} Python, C/C++, Java, JavaScript, SQL, Ada.
  \item \textbf{ML \& LLMs:} PyTorch, Transformers, Hugging Face, Slime (RL framework), LLM post-training (SFT, DPO, RLHF), harness engineering, evaluation, data curation.
  \item \textbf{Program Analysis \& Systems:} LLVM, Z3, static \& dynamic analysis, Linux kernel, fuzzing.
  \item \textbf{Infrastructure:} Linux, Git, Docker, Azure, AWS, distributed training.
  \end{itemize}
\end{minipage}



\end{document}
